% Figaro Paper B3: The Subordination Variable
% Why a Single Coordination Substrate Admits Both the Associated-Producer
% and the Frictionless-Market Reading
% (Third paper in the political-economy trilogy: B1 institutional economics,
% B2 Gramscian hegemony, B3 conditionality of the left/right partition.)
% Preprint, April 2026

\documentclass[11pt,a4paper]{article}

\usepackage[utf8]{inputenc}
\usepackage[T1]{fontenc}
\usepackage{amsmath,amssymb,amsthm}
\usepackage{mathtools}
\usepackage{booktabs}
\usepackage{graphicx}
\usepackage{hyperref}
\usepackage{cleveref}
\usepackage{enumitem}
\usepackage{xcolor}
\usepackage[margin=1in]{geometry}
\usepackage{natbib}
\bibliographystyle{plainnat}

\theoremstyle{remark}
\newtheorem{remark}{Remark}[section]

\newcommand{\figaro}{\textsc{Figaro}}
\newcommand{\fig}{\textsc{Fig}}

\title{%
  The Subordination Variable:\\
  Why a Single Coordination Substrate Admits\\
  Both the Associated-Producer and the Frictionless-Market Reading%
}

\author{
  Alessandro Daliana\thanks{Corresponding author. \texttt{adaliana@figaro.org}}
}

\date{April 2026}

\begin{document}

\maketitle

% ════════════════════════════════════════════════════════════════════
\begin{abstract}
The political-economy traditions usually grouped under the labels
``socialist'' and ``capitalist'' do not agree on much, but they
agree on what divides them. Both traditions identify the wage
relation as the central institutional fact of the modern economy:
one tradition reads it as the locus of value extraction, the other
as the engine of voluntary coordination at scale. In each reading,
the load-bearing variable is the same --- the transfer of residual
control over a worker's conduct to an owner of the means of
production, in exchange for a wage --- and each tradition's
normative conclusions follow from how that transfer is theorized.

This paper observes that the bonded settlement primitive of
the mechanism paper, examined institutionally
in the institutional-economics paper and politically
in the political-economy paper, removes the residual-control
transfer from the set of variables that need to be specified at all.
A wallet bonded to perform an obligation is simultaneously the means
of production, the producer, and the contracting party. There is no
employer to whom residual control is transferred; there is also no
intermediary extracting coordination rent on the transfer's behalf.
Subordination ceases to be a variable, and coordination rent ceases
to be a recurring extraction.

The consequence is that the same artifact admits both readings. A
reader trained in the cooperative tradition can describe the
primitive as the architectural realization of associated production
in Marx's sense (\citealt{marx1867capital}, ch.~1, §4): producers
hold the means of production, coordinate without capital mediation,
and retain the value their labor creates. A reader trained in the
classical-liberal tradition can describe the same primitive as
voluntary exchange between owners stripped of the coordination
rents and authority frictions that have historically attended it
(\citealt{hayek1945knowledge}; \citealt{friedman1962capitalism}).
Both descriptions are accurate. Neither is a stretch.

The substantive claim of this paper is not that the primitive is
``Marxist'' or ``frictionless capitalist''---both descriptions are
category errors made under the assumption that one of them must be
the right answer---but that the partition between the two
traditions was conditional on the subordination variable being
load-bearing in the first place. When the variable is removed,
the partition does not vanish; it relocates. What was a
substrate-level disagreement about the architecture of production
becomes a graph-level disagreement about distribution, public
goods, and the social composition of the assemblies that the
substrate admits. That relocation is the real political-economy
content of the kernel's ideological agnosticism, and it is the
subject of this paper.
\end{abstract}

\medskip
\noindent\textbf{Keywords:}
political economy, subordination, coordination rent, associated
production, classical liberalism, frame collapse, conditionality of
the left/right partition

% ════════════════════════════════════════════════════════════════════
\section{Introduction}
\label{sec:intro}

This paper sits between political philosophy and the institutional
analysis of a particular primitive. It addresses a puzzle that
arose in conversation around the bonded settlement primitive
of the mechanism paper and that proved durable enough to
warrant explicit treatment.

When the primitive is described to a reader trained in the
cooperative or socialist tradition, the natural reading is that
each wallet bonded to perform an obligation \emph{is} a means of
production owned by the actor it represents; that production
proceeds without a wage relation, since there is no party in the
position of employer; that coordination occurs without capital
mediation, since the bonded commitment itself does the
enforcement work that a firm would otherwise do; and that the
value the labor creates is retained by the labor, since there
is no surplus-extraction layer to siphon it. The reading is
recognizable: it is the architectural realization, at protocol
scale, of what \citet{marx1867capital} called the association of
free producers.

When the same primitive is described to a reader trained in the
classical-liberal tradition, an equally natural reading is that
the primitive is voluntary exchange between owners, stripped of the
coordination rents and authority frictions that have historically
attended it. There is no platform extracting a take rate. There
is no firm imposing internal hierarchy. There is no sovereign
gating cross-border settlement. What remains is bilateral
agreement between consenting parties, secured by collateral that
each party voluntarily posts, executed by a public algorithm with
no admin and no fee. The reading is also recognizable: it is the
architectural realization of what
\citet{hayek1945knowledge,friedman1962capitalism} described as
voluntary exchange under decentralized price information, with the
historical impediments to that exchange systematically removed.

The two readings are usually thought to be opposed. They are
classified along an axis that political philosophy has spent two
centuries refining. They occupy distinct vocabularies, name
distinct historical lineages, and recommend distinct institutional
prescriptions. That a single artifact admits both readings, with
neither requiring metaphorical license, is the puzzle this paper
addresses.

\paragraph{The thesis.}
The thesis is not that the primitive is Marxist, capitalist, or
some synthesis or third way. The thesis is structural: the
partition between the two traditions has historically been
organized around two variables --- the subordination relation
between worker and employer, and the coordination rent extracted by
intermediaries that the parties cannot transact without. The
primitive removes both of these variables from the substrate of
exchange. The partition does not survive the removal in its
familiar form, because the two traditions had been arguing about
the same thing using opposed vocabularies, and the thing they had
been arguing about was no longer where they had located it.

\paragraph{What this paper is and is not.}
This paper is a short essay in the conceptual conditions of
political economy. It does not extend the underlying mechanism,
the institutional argument, or the political-economy diagnostic
developed in the companion papers. It takes those contributions
as given and asks what their joint implication is for the
partition that has organized political-economic disagreement
since at least the mid-nineteenth century.

The paper takes no position on which tradition has it right about
distribution, public goods, externalities, or the social
composition of the assemblies that the substrate admits. Those
disagreements are real, and they survive the substrate change.
The paper's claim is narrower: the part of the partition that
appeared to be about the architecture of production was
conditional on the architecture being one in which subordination
and coordination rent were unavoidable. When they become
avoidable, the disagreement does not disappear --- it moves up the
stack. Most of what political economy has been arguing about, on
this reading, was a graph-level question being conducted in
substrate-level vocabulary because the substrate was not yet
disaggregable.

\paragraph{Position relative to the corpus.}
This paper is the third in a political-economy trilogy.
Paper~B1 is the institutional-economics
argument: it shows that the firm's coordination function loses its
cost advantage when bilateral enforcement becomes priceable as a
fixed lockup. Paper~B2 is the
political-economy diagnostic: it identifies the configuration of
firms, platforms, and corridors as a hegemonic arrangement in
Gramsci's sense, contingent in principle and presupposed in
practice. The present paper completes the triangle by examining
what the political-economic stakes are once the variables that
sustained the partition are themselves under question. B1 names
the cost mechanism, B2 names the diagnostic register, B3 names the
partition that both presupposed.

% ════════════════════════════════════════════════════════════════════
\section{The Subordination Variable as Hinge of the Partition}
\label{sec:hinge}

The argument of this section is that two traditions usually
treated as opposed --- the cooperative-Marxian and the
classical-liberal --- agree on the centrality of subordination,
disagree about its meaning, and locate the partition between them
at this disagreement. We then observe that both traditions take
subordination's structural necessity as given, and locate their
prescriptive divergence on top of that shared premise.

\paragraph{The Marxian reading of subordination.}
\citet{marx1867capital} reads the wage relation as the moment at
which a worker's labor capacity is alienated to an owner of the
means of production. The contractual form is voluntary; the
substantive content is the surrender of residual control over the
worker's productive activity to a party whose interests in that
activity are not the worker's. Surplus value is the difference
between what the labor produces and what the wage compensates;
the firm is the institutional container in which the surplus is
realized. The worker, on this reading, cannot retain the full
product of their labor because the means of production is not
theirs --- which is to say, because the architectural condition
of production places means and producer in different hands.
Marx's normative conclusion follows: the way to retain the
product of labor is to reunite the producer with the means of
production, which Marx calls the ``association of free producers''
working with means of production held in common
(\citealt{marx1867capital}, ch.~1, §4). The analytical
reconstruction of this argument by \citet{cohen1978karl} is the
reference we use here for the structural shape of the claim,
without taking a position on the broader historical-materialist
program in which Marx situated it.

\paragraph{The classical-liberal reading of subordination.}
The classical-liberal tradition reads the wage relation as
voluntary exchange between owners of distinct factors of
production: the worker owns labor, the firm owns capital, and the
wage is the price at which the two are combined. The transfer of
residual control to the employer is not an alienation but an
efficient response to the fact that productive activity must be
coordinated and that coordination is costly. \citet{coase1937}
gave the canonical formulation: firms exist because the cost of
discovering prices, negotiating terms, and enforcing contracts in
open markets exceeds the cost of internal hierarchy.
\citet{williamson1985,hart1995} refined the analysis;
\citet{north1990institutions} generalized the move to the level
of institutions broadly construed, reading the institutional
landscape itself as a transaction-cost-minimizing response to
the conditions of exchange. Across these refinements the central
premise is retained: hierarchy and intermediation are efficient
responses to coordination cost. The normative conclusion is that
the wage relation is welfare-improving relative to the
alternative of doing all coordination through spot contracting,
since spot contracting at scale is prohibitively expensive when
enforcement is unreliable.

\paragraph{The shared premise.}
The two traditions disagree about how to interpret subordination
--- exploitation in the first reading, efficiency in the second
--- but agree on the structural fact that some form of it is
required for production to occur at scale under existing
arrangements. Marx does not claim that production can occur
without coordination; he claims that the existing form of
coordination expropriates the producer.
\citet{coase1937,williamson1985,hart1995} do not claim that
coordination is costless; they claim that hierarchy is the
efficient response to its costliness. The disagreement is real,
but it is conducted on a shared terrain.

The shared terrain has a structural feature both traditions take
as given: the only available coordination architectures for
production at scale are (a) the firm, which internalizes
coordination at the cost of subordination, or (b) the spot market,
which avoids subordination at the cost of coordination breakdown.
Both readings of subordination --- as exploitation or as
efficient response to coordination cost --- presuppose this
binary. The cooperative-Marxian tradition's prescriptive program
(reunite producer with means) is hard precisely because it
requires inventing a coordination architecture that the binary
denies the possibility of. The classical-liberal tradition's
descriptive equanimity about the wage relation is sustained
precisely because no such architecture has been available, so the
choice between (a) and (b) genuinely is the relevant choice.

\paragraph{The coordination-rent variable.}
A second variable runs alongside subordination and reinforces the
binary. When parties who have not subordinated themselves to a
common authority must nonetheless transact, the historical
solution has been intermediation: a third party that holds
custody, vouches for performance, or arbitrates dispute, and that
extracts a recurring rent for performing the function.
\citet{srnicek2017platform} traces the contemporary form
(platform capitalism); the form predates platforms by centuries
in the guise of merchant houses, factor agents, clearinghouses,
and the transactional infrastructure of modern finance. Here too,
the cooperative-Marxian tradition reads the rent as expropriation
and the classical-liberal tradition reads it as price for genuine
service; here too, the disagreement is real and the shared
premise is that the rent corresponds to a function that someone
must perform.

\begin{remark}[On reading subordination as a variable]
Treating subordination and coordination rent as ``variables''
that can in principle be set to zero is itself a move that
neither tradition has historically been able to make, because
neither tradition has had access to a coordination architecture
in which they are not load-bearing. The move becomes available
once the architecture changes. The argument of this paper is not
that prior analyses were wrong on their own terms; it is that
their terms held one variable fixed that is no longer fixed.
\end{remark}

% ════════════════════════════════════════════════════════════════════
\section{What the Primitive Does to Both Variables}
\label{sec:what-primitive-does}

We now turn from the partition's structure to what the bonded
settlement primitive of the mechanism paper does to the
two variables on which the partition depends. The mechanism is
established in the companion paper; we summarize only what is
necessary for the present argument.

\paragraph{Asymmetric bonding.}
The primitive runs two mechanisms that compose. The first,
asymmetric bonding, requires the buyer to lock $2P$ in custody
(twice the contract price) and the seller to lock $2G$ in custody
(twice the cumulative seller value the seller is contributing).
Total custody during execution is $2P + 2G$. The bonding ratios
make cooperation the unique strategy profile surviving iterated
elimination of weakly dominated strategies for both parties at
every depth of a multi-party process tree
(Theorem 4.3 of the mechanism paper). Coordination
between the parties is induced by the bond structure itself; no
authority is required to enforce it.

\paragraph{Buyer dominance with atomic resolution.}
The second mechanism, buyer dominance, restricts the resolution
authority to the root buyer and requires that all orders in a
process settle together or not at all. This induces a
weakest-link subgame among co-sellers
(Proposition 5.8 of the mechanism paper) that propagates
cooperation pressure through the process tree without explicit
inter-seller communication. The two mechanisms together produce
the result that an N-party tree of bilateral commitments is
equilibrium-secured from a single buyer signature, with no
intermediary in the position of arbitrator, custodian, or
adjudicator other than the kernel itself --- which holds no
discretion and extracts no rent.

\paragraph{Wallet collapse.}
The companion labor-law paper
establishes the institutional consequence: each wallet
participating in a process is simultaneously the legal property
being transacted (the \emph{res}) and the legal party doing the
transacting (the \emph{persona}). The Roman partition between
property and party that has organized commercial law for two
millennia collapses at the wallet boundary, because the wallet
controls itself: it is what it controls and it controls what it
is. Productive assets become legal subjects; the worker's
productive capacity is no longer a thing rented to someone else
but a thing that acts in its own legal name.

\paragraph{Effect on subordination.}
Subordination, as both traditions have used the term, is the
transfer of residual control over a worker's conduct to an
employer in exchange for a wage. The transfer is not strictly
required by production at scale; it has been required because
prior coordination architectures could not enforce bilateral
agreements between non-subordinated parties at low enough cost.
The bonded primitive enforces such agreements at fixed lockup
cost, with no recurring extraction. The result is that residual
control over a wallet's conduct stays with the wallet's holder.
There is no party to whom control is transferred. Production
proceeds; the wage relation does not.

This is not an empirical claim that the wage relation will
disappear. The wage relation will persist wherever an actor
prefers to operate as an employee --- the substrate is
permissive, not prescriptive. The claim is structural: the wage
relation ceases to be a precondition of coordinated production
at scale. Workers can opt into it, opt out of it, or operate as
their own counterparties. The substrate does not require any of
these choices.

\paragraph{Effect on coordination rent.}
Coordination rent, as both traditions have used the term, is
recurring extraction by an intermediary in exchange for
performing a coordination function the parties cannot perform
themselves. The primitive performs the coordination function ---
custody, enforcement, atomic resolution --- as a public algorithm
with no fee, no admin, and no privileged caller. The bond posted
by the parties is fully reclaimable on cooperation; nothing is
extracted as rent. Intermediaries can still exist on the
substrate (an operator may charge a service fee for coordinating
a process, an aggregator may charge for discovery), but the
substrate does not require them and does not privilege them.
Intermediation becomes a service on the graph, not a tax on the
substrate.

\begin{remark}[Both variables, jointly]
A reader might object that subordination and coordination rent
are the same variable expressed at different scales: the firm
internalizes both, the platform externalizes the second. This is
correct as a historical observation but understates the
generality of the move. The primitive does not relocate the
variables; it removes the function they performed. The
enforcement work formerly done by the firm's authority structure
and the platform's intermediation structure is now done by the
bond structure of the commitment itself, with no party in the
position to extract the work's value as recurring revenue. The
function is performed; no rent attaches.
\end{remark}

% ════════════════════════════════════════════════════════════════════
\section{Two Readings of One Artifact}
\label{sec:two-readings}

We now state the readings that this paper claims are jointly
correct.

\subsection{The Associated-Producer Reading}
\label{sec:associated-producer}

A reader in the cooperative-Marxian tradition encountering the
primitive can describe it as follows. Each wallet bonded to
perform an obligation is a means of production owned by the actor
it represents. The wallet is not a token of identity granted by
an employer; it is a token of self-ownership cryptographically
established and unilaterally controlled. The wallet acts on its
own behalf; it is not subordinated to any other party. Production
proceeds through bilateral commitments between such wallets,
secured by collateral each party posts voluntarily. No party
extracts surplus from another's labor; the bond is reclaimable on
cooperation, and the value the labor creates flows directly to the
labor's wallet at resolution.

The architectural realization of this is, in Marx's own
vocabulary, the association of free producers working with means
of production they hold and coordinating their labor without
capital mediation. The associated producers do not need to
nationalize the means of production, because the means are
already distributed at the wallet boundary. They do not need to
abolish the wage relation, because the wage relation has no
substrate to reproduce on. They do not need to coordinate
production through a planning bureau, because the coordination
function is performed by the bonded commitment itself.

The reading does not require unusual interpretive license. It
follows from the protocol's own behavior. A reader in this
tradition may judge the substrate as the realization of long-held
prescriptive aims and may turn to the question of what graphs
should be composed on it (cooperative-bonding patterns,
contribution-weighted resolution, mutual-aid assemblies). That
question is graph-level, and the substrate admits it.

\subsection{The Frictionless-Market Reading}
\label{sec:frictionless-market}

A reader in the classical-liberal tradition encountering the
same primitive can describe it as follows. Voluntary exchange
between owners is the foundational institution of welfare-improving
production. Historically, voluntary exchange has been impeded by
two friction sources: enforcement (parties who do not trust each
other cannot transact) and intermediation (parties who must use
intermediaries to enforce pay rent for the privilege). Each
friction has been the object of sustained institutional effort:
courts, contract law, payment rails, and platforms have all been
attempts to reduce the first friction at the cost of introducing
or compounding the second.

The bonded primitive removes both frictions in a single move. The
enforcement function is performed by the bond structure of the
commitment, requiring no third party. The intermediation function
is therefore unnecessary, eliminating the rent it extracts.
Voluntary exchange between owners proceeds at substantially lower
friction than has been historically available. The substrate is
the architectural realization of what
\citet{hayek1945knowledge,friedman1962capitalism} described as
the welfare case for voluntary exchange, with the impediments
that have historically prevented it from operating cleanly
systematically removed.

The reading does not require unusual interpretive license either.
It also follows from the protocol's own behavior. A reader in
this tradition may judge the substrate as the realization of
long-held prescriptive aims and may turn to the question of what
graphs should be composed on it (Dutch-auction price discovery,
operator competition, transparent reputation). That question is
also graph-level, and the substrate also admits it.

\subsection{Why Both Readings Are Correct}
\label{sec:both-correct}

The two readings describe the same artifact in different
vocabularies. Each tradition's vocabulary picks out features of
the artifact that the tradition's prior preoccupations have made
salient. The cooperative-Marxian vocabulary picks out the absence
of the wage relation, the producer's ownership of the means of
production, and the absence of surplus extraction. The
classical-liberal vocabulary picks out the voluntary character of
the exchange, the absence of intermediation rent, and the bilateral
form of the commitment. These are descriptions of the same
features under different aspects.

That both readings are correct is not a synthesis. It does not
mean the two traditions secretly agreed all along. It means the
artifact has features that satisfy both descriptions
simultaneously, because the features that the descriptions pick
out are the features that absence-of-subordination and
absence-of-coordination-rent jointly produce, and because the
primitive jointly produces them. The joint description has the
shape: \emph{voluntary exchange between owners of means of
production, none of whom occupies the position of employer of
another, and none of whom is required to pay an intermediary for
the exchange to be enforced.} A cooperative-Marxian reader will
emphasize the first half of that sentence; a classical-liberal
reader will emphasize the second; both halves are true of the
same artifact at the same time.

\begin{remark}[On what is not collapsing]
The two traditions disagree about much more than subordination
and coordination rent. They disagree about distribution, public
goods, externalities, the legitimacy of inheritance, the
appropriate scope of collective decision, and the moral status of
inequality. These disagreements are not addressed by the
substrate change and do not collapse. Section~\ref{sec:relocation}
returns to where they continue to operate. The claim of this
section is narrower: the architectural disagreement that has
been entangled with the distributional disagreement disentangles
when the architectural variable is removed.
\end{remark}

% ════════════════════════════════════════════════════════════════════
\section{Where the Partition Relocates}
\label{sec:relocation}

The partition between the two traditions does not disappear when
the substrate is changed. It relocates. We sketch the relocation
in three places, not exhaustively.

\paragraph{Distribution.}
The substrate determines that bonded commitments settle to the
wallets that posted them; it does not determine how wallets
should be initially endowed, how returns should be shared inside
multi-party assemblies, or how rents arising from network
position should flow. These remain open questions. A graph that
weights resolution payouts by some contribution metric expresses
one position; a graph that distributes uniformly expresses
another; a graph that flows surplus to a treasury for collective
allocation expresses a third. The substrate admits all three.
The contemporary cooperative-tradition literature on this kind
of question --- \citet{wright2010envisioning} on ``real
utopias'' is one anchor --- maps cleanly onto graph-tier
composition once the substrate is no longer the constraint.
The disagreement about which is appropriate is not a substrate
disagreement; it is a graph disagreement. It is not adjudicable
by appeal to the protocol's behavior, because the protocol does
not adjudicate it.

\paragraph{Public goods and externalities.}
The substrate does not produce non-excludable goods, does not
internalize externalities by default, and does not arrange the
collective action by which under-provision of public goods is
addressed. These are first-order questions for any political
economy and are not resolved by the substrate change. A graph
that requires GHG attestation and bonds carbon-removal
commitments expresses one approach; a graph that ignores
externalities expresses another. The substrate admits both. The
question of which approach is right --- whether by efficiency,
justice, or some other criterion --- is a graph-level question
that the substrate does not answer.

\paragraph{The social composition of assemblies.}
The substrate is permissive about who participates. It does not
require a worker, an owner, an artificial intelligence, a state
agency, a cooperative, or a foundation to participate or
abstain; any wallet may compose with any other wallet. The
question of what social composition produces a desirable
political economy --- whether assemblies should privilege
workers, displace platform operators, exclude predatory actors,
admit AI agents on equal footing, route around state authority,
or none of these --- is again a graph-level question. The
substrate admits answers across the spectrum.

\begin{remark}[The kernel is agnostic; the graph is the politics]
The companion political-economy paper (§5) closes with
the formulation: the kernel is ideologically agnostic; the graph
is the politics. The present paper extends that formulation in a
specific way. Once the partition between the two traditions has
been disentangled from the substrate, the disagreement that
remains is not eliminated --- it is concentrated at the graph
tier. The graph is the politics not because the substrate
suppresses politics, but because the substrate frees politics
from being conducted in vocabulary inherited from a particular
configuration of the substrate. The graph is where political
economy actually happens once the architectural variable is no
longer doing its conventional work.
\end{remark}

% ════════════════════════════════════════════════════════════════════
\section{What This Paper Is Not}
\label{sec:not}

The argument is liable to several misreadings. We address them
explicitly.

\paragraph{Not Marxist endorsement.}
The associated-producer reading is offered as a description of
how a reader in that tradition can accurately characterize the
primitive. It is not offered as the protocol's commitment, as
the author's normative endorsement, or as a claim that the
primitive vindicates the Marxian tradition's prescriptive
program. Marx's program included claims about distribution, the
state, class struggle, and historical materialism on which the
substrate takes no position. A reader who concludes from
Section~\ref{sec:associated-producer} that the protocol is
Marxist has read the section as endorsement rather than
description.

\paragraph{Not classical-liberal endorsement.}
Symmetrically, the frictionless-market reading is offered as a
description of how a reader in the classical-liberal tradition
can accurately characterize the primitive. It is not offered as
the protocol's commitment, as the author's normative endorsement,
or as a claim that the primitive vindicates the classical-liberal
prescriptive program. The classical-liberal tradition includes
claims about minimal state, property rights as foundational, and
the moral status of distributional outcomes on which the
substrate takes no position. A reader who concludes from
Section~\ref{sec:frictionless-market} that the protocol is
libertarian has read the section as endorsement rather than
description.

\paragraph{Not synthesis or third way.}
The paper does not propose that the cooperative-Marxian and
classical-liberal traditions agreed all along, that their
disagreement is illusory, or that the primitive is some
hegelian dialectical resolution. The disagreement was real on
its own terms, the variables it disagreed about are real, and
the residual disagreements (Section~\ref{sec:relocation}) remain
real after the substrate change. The paper claims that one part
of the disagreement was conditional on a substrate that admitted
no architectural alternative, and that part disentangles when an
alternative becomes available. The other parts persist.

\paragraph{Not post-political.}
The paper does not claim that political economy is over, that
ideology has ended, or that the substrate's neutrality settles
the questions political economy has historically asked. Quite
the opposite: the substrate's neutrality is the precondition
under which the political-economic questions become legible as
graph-level questions rather than as substrate-level
inevitabilities. The political work that political economy has
always done --- about distribution, public goods, the social
composition of cooperation --- continues, on substrate now
suited to its actual subject matter rather than entangled with
prior architectural constraints. There is more political-economy
work to do, not less.

\paragraph{Not a prediction of adoption.}
Finally, the paper does not predict that either tradition will
recognize the substrate as the realization of its program, that
the primitive will displace incumbent arrangements, or that
political-economic discourse will reorganize itself around the
graph tier in any specific timeframe. The argument is conceptual:
under the supposition that the primitive is widely adopted, the
two traditions' substrate-level disagreement disentangles in the
manner described. Whether adoption occurs is the empirical
question of the political-economy paper, not of this paper.

% ════════════════════════════════════════════════════════════════════
\section{Conclusion}
\label{sec:conclusion}

This paper has argued that the partition between the
cooperative-Marxian and classical-liberal traditions has
historically been organized around two variables --- the
subordination relation and the coordination rent --- both of
which the bonded settlement primitive
of the mechanism paper removes from the substrate of
exchange. The same artifact therefore admits the
associated-producer reading from the first tradition and the
frictionless-market reading from the second, with neither
reading requiring metaphorical license. The disagreement between
the two traditions does not disappear; it relocates from the
substrate to the graph, where the questions it has always been
about (distribution, public goods, social composition of
assemblies) continue to operate.

The substantive observation is that what appeared to be a
foundational political-economic partition was conditional on a
particular configuration of the substrate. The two traditions
were arguing about the same thing using opposed vocabularies,
and the thing they were arguing about was held in place by the
substrate's prior limitations. When the substrate changes, the
partition does not survive in its familiar form. This is not a
victory for either side; it is a relocation of where the genuine
disagreement was always operating.

We close with two observations of method, mirroring the closing
of the companion political-economy paper (§5). First,
the kernel is ideologically agnostic; the graph is the politics
--- and what this paper has added is that the graph is the
politics in a stronger sense than the companion paper stated.
The graph is not just where ideology is expressed once the
substrate is neutral; it is where ideology has always been
operating, partially obscured by the substrate's prior
non-neutrality. The substrate's becoming neutral is what makes
the actual location of ideological disagreement visible.
Second, the disentangling described here is recursive. Future
substrates may relocate further variables from architecture to
choice, and the partition that this paper has described
relocating may itself be relocated again. There is no terminal
relocation; there is only the next variable whose conditionality
we have not yet learned to see.

\section*{Acknowledgments}
The argument of this paper began as an intuition --- that the
bonded primitive admitted both a Marxian and a frictionless-capital
reading of the same protocol --- voiced in conversation. The
intuition was correct as stated; this paper is the unpacking of
why both readings are correct simultaneously, what the joint
correctness implies for the partition the readings come from, and
what survives the relocation. I thank the interlocutor whose
intuition prompted the unpacking, and the readers of earlier
drafts who pressed me to separate the descriptive claim (both
readings are correct) from the normative claim that it would have
been a category error to make (the protocol endorses one reading).


% ════════════════════════════════════════════════════════════════════
% References
% ════════════════════════════════════════════════════════════════════

\begin{thebibliography}{99}

\bibitem[Coase(1937)]{coase1937}
Coase, R.~H.
\newblock The Nature of the Firm.
\newblock \emph{Economica}, 4(16):386--405, 1937.

\bibitem[Cohen(1978)]{cohen1978karl}
Cohen, G.~A.
\newblock \emph{Karl Marx's Theory of History: A Defence}.
\newblock Princeton University Press, Princeton, NJ, 1978.

\bibitem[Friedman(1962)]{friedman1962capitalism}
Friedman, M.
\newblock \emph{Capitalism and Freedom}.
\newblock University of Chicago Press, Chicago, 1962.

\bibitem[Hart(1995)]{hart1995}
Hart, O.
\newblock \emph{Firms, Contracts, and Financial Structure}.
\newblock Oxford University Press, Oxford, 1995.

\bibitem[Hayek(1945)]{hayek1945knowledge}
Hayek, F.~A.
\newblock The Use of Knowledge in Society.
\newblock \emph{American Economic Review}, 35(4):519--530, 1945.

\bibitem[Marx(1867)]{marx1867capital}
Marx, K.
\newblock \emph{Das Kapital, Band I: Der Produktionsprozess des
  Kapitals}.
\newblock Verlag von Otto Meissner, Hamburg, 1867.
\newblock English: \emph{Capital, Volume I}, trans. B.~Fowkes,
  Penguin Classics, 1976.

\bibitem[North(1990)]{north1990institutions}
North, D.~C.
\newblock \emph{Institutions, Institutional Change and Economic
  Performance}.
\newblock Cambridge University Press, Cambridge, 1990.

\bibitem[Srnicek(2017)]{srnicek2017platform}
Srnicek, N.
\newblock \emph{Platform Capitalism}.
\newblock Polity, Cambridge, 2017.

\bibitem[Williamson(1985)]{williamson1985}
Williamson, O.~E.
\newblock \emph{The Economic Institutions of Capitalism}.
\newblock Free Press, New York, 1985.

\bibitem[Wright(2010)]{wright2010envisioning}
Wright, E.~O.
\newblock \emph{Envisioning Real Utopias}.
\newblock Verso, London, 2010.

\end{thebibliography}

\end{document}
